% !PRINTED 445
\weberpage{445}
4. Bedeuten $(k), (k'), \ldots, (k^{(h-1)})$ die sämtlichen Wurzeln der
Gleichung
\[
H_m(u) = 0,
\]
so ist für ein variables $t$
\[
H_m(t) = [t - (k)] [t - (k')] \ldots [t - (k^{(h-1)})],
\]
und diese Funktion hat ganze rationale Koeffizienten. Es ist
also nach (5):
\[
[H_m(t)]^p \equiv H_m(t^p) \quad (\mathrm{mod}\ p),
\]
also, wenn wir $t = (k)$ setzen:
\[
[(k)^p - (k)] [(k)^p - (k')] \ldots [(k)^p - (k^{(h-1)})] \equiv 0 \quad (\mathrm{mod}\ p).
\]
Bedeutet also $\mathfrak P$ ein in $p$ aufgehendes Primideal in $\mathfrak K$, so muß
wenigstens einer der Faktoren der linken Seite durch $\mathfrak P$ teilbar
sein, und daraus folgt eine Kongruenz
\weq{12}{(k)^p \equiv (k') \quad (\mathrm{mod}\ \mathfrak P),}
worin $(k')$ irgend eine der Klasseninvarianten ist, die auch mit
$(k)$ identisch sein kann.

5. Ist zunächst
\[
(D, p) = -1,
\]
so kann $\mathfrak P$ jedenfalls nicht vom ersten Grade sein, weil ja schon
\weq{13}{\begin{aligned}
\sqrt{D}^p &\equiv -\sqrt{D}\\
\sqrt{D}^{p^2} &\equiv \sqrt{D}
\end{aligned} \quad (\mathrm{mod}\ \mathfrak P)}
ist. Ist aber in (12) (nach § 121)
% [sic] The print mixes parenthesized class invariants with bare k in the next three displays.
\[
(k') = f(k, \sqrt{D}), \quad k = f(k', -\sqrt{D}),
\]
so folgt aus (13)
\[
(k')^p \equiv f(k^p, -\sqrt{D}) \equiv f(k', -\sqrt{D}) \equiv k,
\]
\[
(k)^{p^2} \equiv k \quad (\mathrm{mod}\ \mathfrak P),
\]
und mithin ist für alle ganze Zahlen $\xi$ des Körpers $\mathfrak K$
\weq{14}{\xi^{p^2} \equiv \xi \quad (\mathrm{mod}\ \mathfrak P),}
und $\mathfrak P$ ist vom \gesperrt{zweiten Grade}.

6. Ist zweitens
\[
(D, p) = +1,
\]
so gibt es, wie in § 118 bewiesen ist, zwei entgegengesetzte Klassen
$l, l^{-1}$ (die, wenn $l$ zweiseitig ist, miteinander identisch sind), die
durch Formen mit dem ersten Koeffizienten $p$ repräsentiert
werden können, und wenn $(k)$ eine beliebige Klasseninvariante
ist, so ist nach § 120 die Invariantengleichung für den Transformationsgrad $p$:
\weq{15}{F_p(u, v) = 0}
%CONT>>
